\chapter{El modelo insumo-producto de Leontief}
\label{cap:leontief}

\section{Una economía es una red de sectores}

Hasta acá miramos un mercado por vez. Pero una economía real es una \textbf{red de sectores interdependientes}: para producir pan se necesita harina, para producir harina se necesita trigo y energía, y la energía a su vez consume insumos de otros sectores. Lo que es \emph{producto} de un sector es \emph{insumo} de otro. ¿Cómo planificar cuánto debe producir cada sector para abastecer tanto a los demás sectores como a la demanda final de los consumidores?

Wassily \textsc{Leontief} respondió esta pregunta con el \textbf{modelo insumo-producto}, un aporte que le valió el Premio Nobel de Economía en 1973 \cite{leontief}. Es, además, un bellísimo ejemplo de cómo el álgebra lineal —matrices y sistemas de ecuaciones— resuelve un problema económico que de otro modo sería inabordable.

\section{La matriz de coeficientes técnicos}

Supongamos $n$ sectores. La pieza central es la \textbf{matriz de coeficientes técnicos} $A$, cuyo elemento $a_{ij}$ indica \emph{cuántas unidades del sector $i$ se necesitan para producir una unidad del sector $j$}:
\begin{equation}
a_{ij} = \frac{\text{insumo del sector } i \text{ usado por } j}{\text{producción total del sector } j}.
\end{equation}
Cada columna $j$ de $A$ es, entonces, la «receta» del sector $j$: qué insumos y en qué proporción necesita.

Llamemos $\mathbf{x}$ al vector de producción total de cada sector y $\mathbf{d}$ al vector de \textbf{demanda final} (lo que consumen los hogares, el Estado y las exportaciones, por fuera de la cadena productiva). La identidad fundamental dice que \emph{toda la producción se reparte entre insumos intermedios y demanda final}:
\begin{equation}
\mathbf{x} = A\mathbf{x} + \mathbf{d}.
\label{eq:leontief_base}
\end{equation}
El término $A\mathbf{x}$ es lo que los sectores se consumen entre sí; $\mathbf{d}$ es lo que queda para la demanda final.

\section{La solución de Leontief}

La ecuación \eqref{eq:leontief_base} es un sistema lineal con $\mathbf{x}$ como incógnita. Reagrupando:
\begin{equation}
\mathbf{x} - A\mathbf{x} = \mathbf{d} \;\Longrightarrow\; (I - A)\mathbf{x} = \mathbf{d},
\end{equation}
donde $I$ es la matriz identidad. Si la matriz $(I-A)$ es invertible, la producción necesaria para satisfacer cualquier demanda final $\mathbf{d}$ es
\begin{equation}
\boxed{\;\mathbf{x} = (I - A)^{-1}\,\mathbf{d}\;}
\label{eq:leontief_sol}
\end{equation}
A la matriz $L = (I-A)^{-1}$ se la llama \textbf{matriz inversa de Leontief} o \textbf{matriz de requerimientos totales}. Su elemento $\ell_{ij}$ tiene una interpretación preciosa: cuánto debe producir el sector $i$, \emph{en forma directa e indirecta}, por cada unidad de demanda final del sector $j$. Captura los efectos en cadena que a simple vista no se ven.

\begin{ejemplo}[Una economía de dos sectores]
Sean dos sectores, Agro (1) e Industria (2), con matriz de coeficientes técnicos
\[
A = \begin{pmatrix} 0.2 & 0.3 \\ 0.4 & 0.1 \end{pmatrix},
\qquad
\mathbf{d} = \begin{pmatrix} 100 \\ 80 \end{pmatrix}.
\]
Entonces
\[
I - A = \begin{pmatrix} 0.8 & -0.3 \\ -0.4 & 0.9 \end{pmatrix}.
\]
El determinante es $\det(I-A) = (0.8)(0.9) - (-0.3)(-0.4) = 0.72 - 0.12 = 0.60$, y la inversa
\[
(I-A)^{-1} = \frac{1}{0.60}\begin{pmatrix} 0.9 & 0.3 \\ 0.4 & 0.8 \end{pmatrix}
= \begin{pmatrix} 1.5 & 0.5 \\ 0.667 & 1.333 \end{pmatrix}.
\]
La producción requerida es
\[
\mathbf{x} = (I-A)^{-1}\mathbf{d}
= \begin{pmatrix} 1.5(100)+0.5(80) \\ 0.667(100)+1.333(80) \end{pmatrix}
= \begin{pmatrix} 190 \\ 173.3 \end{pmatrix}.
\]
Es decir: para entregar 100 unidades de Agro y 80 de Industria a la demanda final, los sectores deben producir en total \(190\) y \(173.3\) unidades respectivamente. El excedente sobre la demanda final (90 y 93.3) es lo que se consume \emph{dentro} de la cadena productiva.
\end{ejemplo}

\begin{ideabox}
\textbf{Lectura de gestión.} El modelo de Leontief es la base de la planificación de la producción y del cálculo de \emph{multiplicadores} económicos: si la demanda de un sector crece, ¿cuánto debe crecer toda la economía para sostenerla? La matriz $(I-A)^{-1}$ contiene esa respuesta, efectos indirectos incluidos.
\end{ideabox}

\begin{labbox}
Acá el álgebra a mano se vuelve impracticable apenas hay más de dos o tres sectores: invertir una matriz $10\times 10$ es imposible con lápiz y papel pero instantáneo para una computadora. En el laboratorio se arma $A$ con \texttt{numpy}, se calcula \texttt{I - A} y la producción con \texttt{numpy.linalg.solve(I - A, d)} (numéricamente más estable que invertir y multiplicar). Es el ejemplo más claro de por qué necesitamos la máquina: la teoría cabe en una fórmula, pero el cálculo no cabe en una hoja.
\end{labbox}

\section{Una nota sobre el solapamiento con Álgebra}

Si ya cursaste Álgebra, parte de este capítulo te va a resultar familiar: la inversión de matrices y los sistemas lineales son contenido de esa materia. La diferencia —y el aporte de este curso— es la \textbf{interpretación económica}: acá una matriz no es un arreglo abstracto de números sino la radiografía de cómo una economía se produce a sí misma. Ese cambio de mirada, de la mecánica del cálculo al significado, es justamente lo que buscamos.

\section{Síntesis del capítulo}

El modelo de Leontief generaliza la idea de equilibrio a una economía con muchos sectores interdependientes, y la resuelve con la inversa $(I-A)^{-1}$. Con esto cerramos el bloque «estático» del curso. A partir del próximo capítulo incorporamos el \emph{cambio}: cómo medir variaciones con derivadas y por qué el pensamiento marginal está en el centro de toda decisión económica.
